\section{Limitations and Ethics}

The current study has four main limitations.
First, the benchmark is entirely synthetic even though the accounting
mechanics are realistic.
Synthetic control is valuable for exact scoring, but it cannot cover
the full distribution of language, ambiguity, and disclosure
conventions found in live corporate reporting.

Second, the current main-paper results are concentrated on the
Controlled-100 split.
That design is appropriate for paired comparison across expensive
systems, but a submission-ready benchmark should also report 500-problem
results for smaller open-weight models so that the strongest empirical
claims rest on 500 or more instances.

Third, the present paper studies the constructor role but not yet the
auditor role.
In practical accounting workflows, a useful system must both construct
candidate statements and audit them for inconsistencies.
Extending \benchmark{} into a companion auditing benchmark is therefore
an important part of the broader project scope.

Fourth, all experiments in this draft rely on end-to-end language-model
generation, which means the model is responsible for both selecting the
correct accounting operations and performing the arithmetic internally.
Because bookkeeping arithmetic is deterministic, an important next step
is an agentic variant in which the model issues typed function calls to
an external calculator or ledger engine while the tool performs the
actual computation.
That comparison would clarify how much of current failure is due to
accounting misclassification versus arithmetic execution.

From an ethics perspective, the benchmark should be used as an
evaluation tool rather than evidence that a model is safe for autonomous
accounting use.
Several systems produce fluent, well-structured outputs while violating
the accounting equation or silently mishandling equity closure.
That failure mode is precisely the kind that could evade casual human
inspection.
